% experiment_comparison.tex
% Section for cross-experiment comparison and key findings

\section{Comparative Analysis and Key Findings}
\label{sec:comparative_analysis}

This section summarizes and compares the results and key insights across all four experiments conducted in this work: MNIST (Exp.\ 1), CIFAR-10 (Exp.\ 2), WikiArt (Exp.\ 3), and CelebA (Exp.\ 4). Each experiment was designed to incrementally increase task difficulty along multiple axes simultaneously: output resolution, dataset diversity, caption richness, and architectural complexity. Together they trace a coherent progression from a minimal proof-of-concept to a high-resolution latent diffusion system.

\subsection{Configuration Comparison}

Table~\ref{tab:all_experiments_comparison} summarizes the key configuration choices across all four experiments.

\begin{table}[htbp]
    \centering
    \caption{Side-by-side configuration comparison of all four experiments.}
    \label{tab:all_experiments_comparison}
    \small
    \begin{tabular}{lcccc}
        \toprule
        \textbf{Aspect} & \textbf{MNIST} & \textbf{CIFAR-10} & \textbf{WikiArt} & \textbf{CelebA} \\
        \midrule
        Resolution & 28$\times$28 & 32$\times$32 & 128$\times$128 & 256$\times$256 \\
        Channels & 1 (grayscale) & 3 (RGB) & 3 (RGB) & 3 (RGB) \\
        Diffusion space & Pixel & Pixel & Pixel & Latent (32$\times$32$\times$4) \\
        Classes/Attributes & 10 digits & 10 objects & 27 styles & 40 binary attributes \\
        Training samples & 60,000 & 50,000 & $\sim$81,000 & $\sim$202,599 \\
        U-Net parameters & $\sim$3M & $\sim$45M & $\sim$150M & $\sim$90M \\
        Additional models & None & None & None & \texttt{stabilityai/sd-vae-ft-mse} \\
        Batch size & 512 & 128 & 64 & 128 \\
        Learning rate & $10^{-3}$ & $10^{-4}$ & $10^{-5}$ & $10^{-5}$ \\
        Training epochs & 20 & 50 & 100 & 500 \\
        Caption type & Class-label template & Class-label template & Class-label template & Attribute-based NLG\footnotemark \\
        Text encoder & \multicolumn{4}{c}{Frozen CLIP (\texttt{openai/clip-vit-base-patch32})} \\
        \bottomrule
    \end{tabular}
\end{table}
\footnotetext{NLG: Natural Language Generation. CelebA captions are programmatically assembled from the 40 binary attributes rather than using a fixed template.}

\subsection{Scaling Trends}

A consistent pattern of deliberate scaling is visible across all experiments. As dataset complexity and target resolution increased, three aspects of the configuration were adjusted in concert: model capacity, learning rate, and training duration.

\textbf{Model capacity.}
The U-Net grew from $\sim$3M parameters for MNIST to $\sim$150M for WikiArt, reflecting the need to model richer feature distributions at higher resolutions. Interestingly, the CelebA model ($\sim$90M parameters) is smaller than WikiArt despite producing higher-resolution outputs (256$\times$256 vs.\ 128$\times$128 pixels). This is possible because CelebA operates in the VAE's 32$\times$32 latent space: the U-Net never processes full-resolution tensors, so it does not need to be sized for 128$\times$128 spatial operations.

\textbf{Learning rate and batch size.}
Learning rate decreased from $10^{-3}$ (MNIST) to $10^{-4}$ (CIFAR-10) to $10^{-5}$ (WikiArt, CelebA), following the standard practice of reducing the step size as model capacity and data complexity grow. Batch size also decreased from 512 (MNIST) to 64 (WikiArt) as the memory footprint of each sample grew with resolution, before recovering to 128 for CelebA because latent-space inputs are compact.

\textbf{Training duration.}
The number of required training epochs grew substantially: 20 (MNIST), 50 (CIFAR-10), 100 (WikiArt), and 500 (CelebA). This reflects both greater dataset diversity and the richer conditional distributions the model must learn. The CelebA model's 500-epoch requirement is especially notable: the 40-attribute binary space induces a very large combinatorial variety of captions, requiring the conditioning pathway to associate many different text descriptions with spatially similar face images.

\subsection{Prompt Conditioning Strategies}

Experiments 1--3 used a fixed-template captioning strategy: a single string template parameterized by the class label (e.g., \texttt{"A handwritten digit 5"}, \texttt{"A photo of a cat"}, \texttt{"A painting in the style of Impressionism"}). This approach has two advantages: it is simple to implement and the vocabulary of conditioning signals is small (10, 10, and 27 unique captions respectively), making it easy for the model to align visual features with prompt tokens.

Experiment 4 replaces fixed templates with a rule-based natural language generation (NLG) system that converts the 40 binary CelebA attributes into diverse, sentence-form captions. Each training sample receives a unique caption assembled from whichever attributes are active, with paraphrase variety, collision resolution, and controlled inclusion of minor features. This approach offers several important advantages over fixed-template conditioning:

\begin{itemize}
    \item \textbf{Richer supervision signal:} Each training caption is distinct, preventing the model from memorizing a small set of global class representations.
    \item \textbf{Fine-grained inference control:} At inference time, any combination of attribute phrases can be assembled into a prompt, providing granular control over gender, hair colour, expression, accessories, and more.
    \item \textbf{Natural language alignment:} Prompts like \texttt{"A photo of a young woman with straight blond hair, smiling, wearing earrings"} exploit CLIP's broad natural-language training, allowing the text encoder to generalize across phrasings not seen during diffusion training.
\end{itemize}

The principal trade-off is increased training complexity: the model must learn to associate hundreds of thousands of unique captions with the underlying visual attributes rather than mapping a handful of fixed strings to global class prototypes.

\subsection{Pixel-Space vs.\ Latent-Space Diffusion}

The first three experiments operate directly in pixel space: the denoising U-Net takes noisy pixel tensors as input and outputs noise predictions in the same space. This is practical at low resolutions (28--128 pixels), but scaling to 256$\times$256 in pixel space would have required processing tensors of 196,608 elements per sample---a 48$\times$ increase over the 4,096-element 32$\times$32$\times$4 latents used in Experiment 4.

Experiment 4 adopts a \emph{latent diffusion} approach inspired by LDMs~\cite{rombach2022high}. A frozen, pretrained VAE (\url{stabilityai/sd-vae-ft-mse}) compresses each 256$\times$256 RGB image to a 32$\times$32$\times$4 latent by an 8$\times$ spatial factor. The U-Net is trained entirely in this latent space and only the VAE decoder is used at inference time to reconstruct the final image. The consequences are significant:

\begin{itemize}
    \item \textbf{Memory efficiency:} GPU memory per sample scales with $(H/8)\times(W/8)$ rather than $H\times W$, making 256$\times$256 generation feasible within the same memory budget as 32$\times$32 pixel-space models.
    \item \textbf{Training stability:} The training loss in latent space exhibited smooth, monotonic convergence, which stands in contrast to the oscillations sometimes observed in pixel-space diffusion at higher resolutions.
    \item \textbf{Frozen model dependency:} The latent diffusion approach introduces a pretrained VAE as a frozen component. This adds a model to load and maintain, and the quality of reconstructed images is upper-bounded by the VAE's reconstruction fidelity.
\end{itemize}

\subsection{Key Takeaways}

\begin{enumerate}
    \item \textbf{A unified architecture scales across four orders of magnitude in dataset complexity.} The same core design --- frozen CLIP text encoder, cross-attention U-Net, DDPM scheduler, CFG --- was applied without fundamental changes from 28$\times$28 grayscale digits to 256$\times$256 celebrity faces, demonstrating the generality of the text-conditioned diffusion paradigm.

    \item \textbf{Latent diffusion is essential for high-resolution generation under practical compute budgets.} Pixel-space diffusion is viable up to $\sim$128$\times$128; beyond that, the memory and compute cost grows quadratically. Offloading spatial compression to a frozen pretrained VAE decouples resolution scaling from U-Net size.

    \item \textbf{Data engineering effort grows with dataset scale.} The WikiArt data-loading bottleneck and the CelebA NLG system both illustrate that, at scale, the engineering work surrounding training data preparation can exceed the effort spent on model architecture design. Efficient data pipelines are as important as model choices for practical training.
\end{enumerate}
